\bibitem{formal-arcs-lan}
Kai-Wen Lan, \emph{Arithmetic compactifications of PEL-type Shimura
varieties}, author's revised thesis, 14 March 2021.
Lemma 4.2.1.6; Theorems 4.2.1.14 and 4.4.16;
Sections 4.3.1 and 4.5; Corollary 4.5.1.8 and its following
affine-chart formula; Proposition 4.5.2.15;
Theorem 4.5.3.10, Proposition 4.5.3.11 and Corollary 4.5.3.12;
Propositions 3.3.1.5 and 3.3.2.2.
\url{https://www.kwlan.org/articles/cpt-PEL-type-thesis-revision.pdf}.
The construction uses the complete pair $(O'[[t]],(t))$ and
the original primitive-line period action.

\bibitem{formal-arcs-ramified}
Kai-Wen Lan, \emph{Compactifications of PEL-type Shimura varieties
in ramified characteristics}, Forum of Mathematics, Sigma 4 (2016),
e1, 98 pp.  Theorem 11.4 and the exact cone condition of
Theorem 7.14.
\url{https://www.kwlan.org/articles/cpt-ram-flat.pdf}.
The normal-base extension theorem is applied to the actual
two-dimensional power-series base, whose monodromy is one ray.

\bibitem{formal-arcs-source}
Levent Alp\"oge, pinned \texttt{jsmorph/rh} source at
\texttt{ecd09ff9b4aaf6c08f2b5ef17dc49ca9648a77de}.
\texttt{UniformisationConclusion}, lines 130--170 and 182--227;
\texttt{DegUnif}, lines 51--76 and 105--178;
\texttt{DegExp}, lines 91--144; and \texttt{GtPtBridgeSt},
lines 90--143.  The certificate records the exact source
interfaces and all normalization signs used in the proof.

\bibitem{formal-arcs-relative-proj}
The Stacks Project, \emph{Relatively ample sheaves},
Tag \texttt{01VJ}.  The canonical relative-Proj open immersion
for a relatively ample invertible sheaf is used to compare
the two actual formal coordinate maps.
\url{https://stacks.math.columbia.edu/tag/01VJ}.
